% ---
% titre: 9P - Lignes et surfaces - Corrigés (réponses seules)
% theme: GM
% chapitre: Lignes et surfaces
% sous_chapitre: 
% niveau: [9e]
% type: RES
%categorie: corrigé
% tags: c-conversion-unite,c-perimetre,c-aire, f-corrige, p-auto-correction
% corrige: false
% source: latex
% ---

% ---- Macros locales à ce fichier (pas dans template_site.tex) ----
\newcommand{\etape}[1]{%
  \vspace{10pt}
  {\Large\bfseries\textcolor{themeColor}{#1}}
  \vspace{4pt}
  \par
}
 
 % Conversion : question = réponse (réponse en gras)
\newcommand{\cv}[4]{#1~#2\ $=$\ \textbf{#3}~#4}
 
% =========================================================
\etape{1 \quad Conversion d'unités}
% =========================================================
 
\header{Problème 60 --- Conversions}
{\small
\renewcommand{\arraystretch}{1.1}
\textbf{A.}\\[2pt]
\begin{tabularx}{\textwidth}{XX}
\cv{13}{m}{1\,300}{cm} & \cv{5,5}{hm}{550}{m} \\
\cv{180}{cm}{1,8}{m} & \cv{0,8}{m}{800}{mm} \\
\cv{5}{cm}{0,5}{dm} & \cv{2,4}{km}{24}{hm} \\
\cv{0,4}{dam}{4\,000}{mm} & \cv{9}{cm}{0,009}{dam} \\
\cv{800}{m}{80}{dam} & \cv{42}{cm}{420}{mm} \\
\cv{1\,300}{mm}{13}{dm} & \cv{0,3}{km}{300}{m} \\
\end{tabularx}
 
\vspace{5pt}
\textbf{D.}\\[2pt]
\begin{tabularx}{\textwidth}{XX}
\cv{832}{m\textsuperscript{2}}{8,32}{dam\textsuperscript{2}} & \cv{0,08}{m\textsuperscript{2}}{800}{cm\textsuperscript{2}} \\
\cv{4\,000}{dm\textsuperscript{2}}{400\,000}{cm\textsuperscript{2}} & \cv{5}{dam\textsuperscript{2}}{0,0005}{km\textsuperscript{2}} \\
\cv{0,7}{km\textsuperscript{2}}{7\,000}{dam\textsuperscript{2}} & \cv{52\,800}{dm\textsuperscript{2}}{0,0528}{hm\textsuperscript{2}} \\
\cv{17,5}{mm\textsuperscript{2}}{0,175}{cm\textsuperscript{2}} & \cv{125,7}{cm\textsuperscript{2}}{12\,570}{mm\textsuperscript{2}} \\
\cv{2\,500\,000}{cm\textsuperscript{2}}{2,5}{dam\textsuperscript{2}} & \cv{748}{mm\textsuperscript{2}}{7,48}{cm\textsuperscript{2}} \\
\cv{1\,767}{hm\textsuperscript{2}}{17,67}{km\textsuperscript{2}} & \cv{0,075}{hm\textsuperscript{2}}{7,5}{dam\textsuperscript{2}} \\
\end{tabularx}
 
\vspace{5pt}
\textbf{E.}\\[2pt]
\begin{tabularx}{\textwidth}{XX}
\cv{0,16}{dm\textsuperscript{2}}{1\,600}{mm\textsuperscript{2}} & \cv{2\,002}{m\textsuperscript{2}}{0,002\,002}{km\textsuperscript{2}} \\
\cv{136,75}{dam\textsuperscript{2}}{136\,750\,000}{cm\textsuperscript{2}} & \cv{72,5}{hm\textsuperscript{2}}{72\,500\,000}{dm\textsuperscript{2}} \\
\cv{5}{mm\textsuperscript{2}}{0,05}{cm\textsuperscript{2}} & \cv{3,14}{m\textsuperscript{2}}{0,0314}{dam\textsuperscript{2}} \\
\cv{0,067}{km\textsuperscript{2}}{6\,700\,000}{dm\textsuperscript{2}} & \cv{0,117}{dam\textsuperscript{2}}{11\,700\,000}{mm\textsuperscript{2}} \\
\cv{450\,000}{hm\textsuperscript{2}}{4\,500}{km\textsuperscript{2}} & \cv{16\,535}{dm\textsuperscript{2}}{0,016\,535}{hm\textsuperscript{2}} \\
\cv{345}{cm\textsuperscript{2}}{0,0345}{m\textsuperscript{2}} & \cv{3\,000}{cm\textsuperscript{2}}{0,3}{m\textsuperscript{2}} \\
\end{tabularx}
 
\vspace{5pt}
\textbf{F.}\\[2pt]
\begin{tabularx}{\textwidth}{XX}
\cv{0,32}{dam\textsuperscript{2}}{3\,200}{dm\textsuperscript{2}} & \cv{9}{dm\textsuperscript{2}}{0,0009}{dam\textsuperscript{2}} \\
\cv{0,0006}{m\textsuperscript{2}}{600}{mm\textsuperscript{2}} & \cv{81}{m\textsuperscript{2}}{0,000\,081}{km\textsuperscript{2}} \\
\cv{4,8}{hm\textsuperscript{2}}{48\,000}{m\textsuperscript{2}} & \cv{0,04}{dm\textsuperscript{2}}{400}{mm\textsuperscript{2}} \\
\cv{575}{dm\textsuperscript{2}}{5,75}{m\textsuperscript{2}} & \cv{1,23}{km\textsuperscript{2}}{1\,230\,000}{m\textsuperscript{2}} \\
\cv{690}{cm\textsuperscript{2}}{69\,000}{mm\textsuperscript{2}} & \cv{8\,960}{cm\textsuperscript{2}}{0,896}{m\textsuperscript{2}} \\
\cv{85\,000}{mm\textsuperscript{2}}{0,000\,85}{dam\textsuperscript{2}} & \cv{748}{dam\textsuperscript{2}}{7,48}{hm\textsuperscript{2}} \\
\end{tabularx}
}
 
\vfill
% =========================================================
\etape{2 \quad Aire et périmètre}
% =========================================================
 
\header{GM26 --- Aires de triangles}
\vspace{-10pt}\begin{tabularx}{\textwidth}{XXX}
$A_{ABC} = 4{,}34$ & $A_{EFG} = 6$ & $A_{IJK} = 4{,}995$ \\
\end{tabularx}
 
\vspace{20pt}
\header{GM27 --- Aires de quadrilatères}
\vspace{-10pt}\begin{tabularx}{\textwidth}{XXX}
$A_{ABCD} = 5{,}76$ & $A_{FGHI} = 6{,}72$ & $A_{JKLM} = 4$ \\
\end{tabularx}
 
\newpage
\header{Problème 92 --- C'est un fameux deux-mâts\,\ldots}
\vspace{-5pt}
{\small
\renewcommand{\arraystretch}{1.3}
\begin{tabularx}{0.45\textwidth}{l|C}
Partie & Aire (cm$^2$) \\
\hline
Coque        & 46,8 \\
Mât gauche   & 4 \\
Mât droit    & 4 \\
Voile gauche & 32 \\
\hline
\end{tabularx}\hfill
\begin{tabularx}{0.45\textwidth}{l|C}
Partie & Aire (cm$^2$) \\
\hline
Voile centre & 36,4 \\
Voile droite & 22,4 \\
Cabine       & 3,2 \\
Drapeau      & 0,64 \\
\hline
\end{tabularx}\hfill
\renewcommand{\arraystretch}{1}

\vspace{5pt}\textbf{Total}: 149,44\,cm$^2$}

 
\vspace{20pt}
\header{GM36 --- Dans le ciel}
\vspace{-10pt}$A = 1\,000$ cm$^2$
 
\vspace{20pt}
\header{GM33 --- Aires de carrés}
\vspace{-10pt}\begin{tabularx}{\textwidth}{XX}
Aire = 274 cm$^2$ & Périmètre = 80 cm \\
\end{tabularx}
 
\vspace{20pt}
\header{GM34 --- Des simples et des composées}
\vspace{-5pt}
{\small
\renewcommand{\arraystretch}{1.3}
\begin{tabularx}{\textwidth}{l|*{5}{C}}
Figure & a) & b) & c) & d) & e) \\
\hline
Périmètre en dm & 82 & 8 & $\approx$ 2,24 & $\approx$ 29,6 & 4 \\
Aire en dm$^2$ & 312 & 2,7 & 0,31 & $\approx$ 65,69 & 0,5025 \\
\end{tabularx}}
 
\vspace{20pt}
\header{GM49 --- À vendre}
\vspace{-10pt}Aire du terrain: 960 m$^2$\\
Prix du terrain: CHF 177\,600.--
 
\vspace{20pt}
\header{GM52 --- Du périmètre à l'aire}
\vspace{-10pt}Aire = 14,7 cm$^2$ (ou 1\,470 mm$^2$)
 
\vspace{20pt}
\header{GM51 --- Un gros mangeur}
\vspace{-10pt}Le côté de la pizza carrée de Joe mesure 30 cm.
 
\newpage
% =========================================================
\etape{3 \quad Exercices \og\,Défi \fg}
% =========================================================
 
\header{GM53 --- Pas si simple qu'il en a l'air}
\vspace{-10pt}$A_{ADC} = 2{,}25$ cm$^2$
 
\vspace{20pt}
\header{GM50 --- Football et mesures}
\vspace{-15pt}\begin{enumerate}[leftmargin=*,align=parleft, label=\alph*)]
\item 8\,436 m$^2$
\item 374 m
\end{enumerate}
 
\vspace{20pt}
\header{Problème 89 --- MAireTH}
\vspace{-5pt}
{\small
\renewcommand{\arraystretch}{1.3}
\begin{tabularx}{0.45\textwidth}{cX|>{\centering\arraybackslash}p{2cm}}
 & Figure & Aire (cm$^2$) \\
\hline
A. & Trapèze rectangle & 8,28 \\
B. & Triangle rectangle & 1,08 \\
C. & Cerf-volant & 6,21 \\
D. & Parallélogramme & 2,295 \\
E. & Grand et petit triangles & 11,415 \\
\hline
\end{tabularx}\hfill
\begin{tabularx}{0.45\textwidth}{cX|>{\centering\arraybackslash}p{2cm}}
 & Figure & Aire (cm$^2$) \\
\hline
F. & Trapèze rectangle & 8,505 \\
G. & Rectangle & 4,06 \\
H. & Triangle rectangle & 5,985 \\
I. & Carré & 2,89 \\
J. & Triangle & 5,13 \\
\hline
\end{tabularx}

\vspace{5pt}\textbf{Surface grise totale}:  55,85\,cm$^2$}